
\documentclass{article}
\usepackage{colortbl}
\usepackage{makecell}
\usepackage{multirow}
\usepackage{supertabular}

\begin{document}

\newcounter{utterance}

\centering \large Interaction Transcript for game `wordle\_withcritic\_noexp', experiment `medium\_frequency\_words\_clue\_with\_critic', episode 19 with reflexion{-}Qwen3.5{-}9B.
\vspace{24pt}

{ \footnotesize  \setcounter{utterance}{1}
\setlength{\tabcolsep}{0pt}
\begin{supertabular}{c@{$\;$}|p{.15\linewidth}@{}p{.15\linewidth}p{.15\linewidth}p{.15\linewidth}p{.15\linewidth}p{.15\linewidth}}
    \# & $\;$A & \multicolumn{4}{c}{Game Master} & $\;\:$B\\
    \hline

    \theutterance \stepcounter{utterance}  
    & & \multicolumn{4}{p{0.6\linewidth}}{
        \cellcolor[rgb]{0.9,0.9,0.9}{
            \makecell[{{p{\linewidth}}}]{
                \texttt{\tiny{[P1$\langle$GM]}}
                \texttt{You have attempted this same task before and did not succeed. These are the plans you wrote after those attempts.} \\
\texttt{Trial \#1: New plan: I failed because I generated a response containing only the "guess:" tag without the required "explanation:" tag, violating the implicit format expected by the Game Master (as seen in Example 2 where both tags were mandatory for a valid turn). Additionally, I ignored the specific clue "tend the hearth" by guessing the exact word "hearth" immediately, which is risky if the target is a synonym or related concept. Next time, I will first construct a response that includes both an "explanation:" tag detailing my reasoning based on the clue and a "guess:" tag containing a valid five{-}letter word that fits the clue but is not the obvious answer (e.g., "ember" or "fire" if they fit the length and context, though "hearth" is 6 letters, so I must choose a 5{-}letter synonym like "stove" or "coals" if valid, or a related 5{-}letter word). I will ensure the output strictly follows the "explanation: {[}text{]}\textbackslash{}nguess: {[}word{]}" format.} \\
\\ 
\\ 
\texttt{You are a language wizard who likes to guess words by using the given rules.} \\
\\ 
\texttt{Welcome to Wordle! You have six attempts to guess the target word, a valid English word of five lowercase letters (a{-}z). Please use the tag "guess:" to provide your guess.} \\
\\ 
\texttt{To help you make an informed guess, you will receive a clue for the word, such as} \\
\texttt{clue: "snowy white"} \\
\\ 
\texttt{Here is an example guess based on the clue: } \\
\texttt{guess: apple} \\
\\ 
\texttt{I will then indicate whether I agree or disagree with your guess, but agreeing with a guess does not confirm its correctness. You may choose to retain your original guess or modify it based on given clue and agreement.} \\
\\ 
\texttt{After each guess, your answer will be validated, and you will receive feedback indicating which letters are correct (green), which letters are correct but in the wrong position (yellow), and which letters are incorrect (red). This feedback can be useful in determining which letters to include or exclude in your next guess.} \\
\\ 
\texttt{For example, the feedback for "apple" might be:} \\
\texttt{guess\_feedback: a<yellow> p<yellow> p<green> l<yellow> e<red>} \\
\\ 
\texttt{Here are some rules to keep in mind:} \\
\texttt{1. Please do not generate any guess feedback or clue on your own} \\
\texttt{2. Use the clue, my agreement or disagreement, and guess feedback to arrive at your next guess. } \\
\\ 
\texttt{Do not generate any other text.} \\
\\ 
\texttt{Let's begin with your first guess.} \\
\\ 
\texttt{clue: tend the hearth} \\
            }
        }
    }
    & \\ \\

    \theutterance \stepcounter{utterance}  
    & \multicolumn{4}{p{0.6\linewidth}}{
        \cellcolor[rgb]{0.9,0.9,0.9}{
            \makecell[{{p{\linewidth}}}]{
                \texttt{\tiny{[P1$\rangle$GM]}}
                \texttt{explanation: The clue "tend the hearth" suggests a word related to maintaining a fire or the area where a fire burns. Since "hearth" is 6 letters, I need a 5{-}letter synonym or related term. "Stove" is a common 5{-}letter word associated with a hearth, but it might be too specific. "Ember" fits the context of a hearth and is 5 letters. However, "tend" implies an action or a role. "Keeper" is 6 letters. "Fire" is 4 letters. "Burn" is 4 letters. "Kindle" is 6 letters. "Warm" is 4 letters. "Heat" is 4 letters. "Coals" is 5 letters and directly relates to what one tends in a hearth. Let's try "coals".} \\
\texttt{guess: coals} \\
            }
        }
    }
    & & \\ \\

    \theutterance \stepcounter{utterance}  
    & & & \multicolumn{2}{p{0.3\linewidth}}{
        \cellcolor[rgb]{0.9,0.9,0.9}{
            \makecell[{{p{\linewidth}}}]{
                \texttt{\tiny{[GM$|$GM]}}
                \texttt{Error: The response should always start with the keyword 'guess:'} \\
            }
        }
    }
    & & \\ \\

    \theutterance \stepcounter{utterance}  
    & & & \multicolumn{2}{p{0.3\linewidth}}{
        \cellcolor[rgb]{0.9,0.9,0.9}{
            \makecell[{{p{\linewidth}}}]{
                \texttt{\tiny{[GM$|$GM]}}
                \texttt{game\_result = ABORT} \\
            }
        }
    }
    & & \\ \\

    \theutterance \stepcounter{utterance}  
    & \multicolumn{4}{p{0.6\linewidth}}{
        \cellcolor[rgb]{0.95,0.95,0.95}{
            \makecell[{{p{\linewidth}}}]{
                \texttt{\tiny{[P2$\rightarrow$P1]}}
                \texttt{Post{-}episode Reflexion (not sent to GM):} \\
\texttt{New plan: I failed because my previous response included an "explanation:" tag, which was not part of the required format for this specific game instance. The rules explicitly stated to use only the "guess:" tag and to generate no other text. Additionally, while "coals" is a valid 5{-}letter word fitting the clue "tend the hearth," the primary cause of the abort was the format violation. Next time, I will strictly output only the "guess: {[}word{]}" line without any preceding explanation or extra text, ensuring the response adheres exactly to the "Do not generate any other text" constraint.} \\
            }
        }
    }
    & & \\ \\

\end{supertabular}
}

\end{document}
